\section{Experiments}

\subsection{Datasets}

To comprehensively evaluate the proposed Unified Dual-Mask Physical Model, we conduct extensive experiments on three distinct light field datasets, carefully selected to cover a wide spectrum of surface reflectance properties and scene complexities. Specifically, our evaluation encompasses Lambertian, Mixed/Urban, and non-Lambertian domains. This diverse selection is essential for validating the robustness of our model, particularly in addressing the inherent limitations of traditional Epipolar Plane Image (EPI) based methods when confronted with non-Lambertian reflections and complex textures.

\textbf{HCI New Dataset.} The HCI New dataset [1] is a widely used synthetic benchmark primarily consisting of Lambertian scenes with complex textures, intricate geometries, and occlusions (e.g., \textit{boxes}, \textit{cotton}, \textit{dino}). Each scene comprises 81 sub-aperture images arranged in a $9 \times 9$ angular grid, with a native spatial resolution of $512 \times 512$. The ground-truth depth maps are rendered using Blender [2] with perfect Lambertian reflectance. We partition the dataset into training, validation, and testing sets at the scene level to ensure strict data isolation. This dataset serves to evaluate the model's fundamental depth estimation capabilities and its performance limits on complex Lambertian textures, where EPI slopes tend to blur due to high-frequency spatial details.

\textbf{UrbanLF-Syn Dataset.} To assess the model's generalization in complex, real-world-like environments, we employ the UrbanLF-Syn dataset [3]. It contains 170 synthetic scenes depicting urban and outdoor environments, characterized by mixed reflectance properties (i.e., a combination of Lambertian and mild non-Lambertian surfaces) and large-scale structures. The angular resolution is configured to $9 \times 9$, and the spatial resolution is uniformly preprocessed to $512 \times 512$. The scenes are divided into training, validation, and testing splits. This dataset is instrumental in evaluating the model's efficacy in mixed domains, where varying illumination and diverse material properties challenge conventional photo-consistency assumptions.

\textbf{non-Lambertian Dataset.} Traditional light field depth estimation methods heavily rely on the Lambertian assumption, leading to severe performance degradation on specular, glossy, or scattering surfaces. To explicitly evaluate our model on such challenging physical properties, we utilize a specialized non-Lambertian dataset [4]. This dataset features scenes with prominent non-Lambertian materials, where the EPI structure is significantly distorted by view-dependent reflections. A critical challenge of this dataset is its extreme data scarcity, containing only 4 scenes for training. The angular and spatial resolutions are maintained at $9 \times 9$ and $512 \times 512$, respectively. The severe lack of training data makes this dataset an ideal testbed to verify whether our proposed dual-mask physical model can learn robust physical priors without overfitting, thereby overcoming the physical failure of EPI assumptions.

\textbf{Data Preprocessing and Sampling Strategy.} 
All datasets are preprocessed to ensure a unified input format for our network. Specifically, sub-aperture images are cropped or resized to a consistent spatial resolution of $512 \times 512$, and a $9 \times 9$ angular grid is extracted. The ground-truth depth maps are normalized to a unified disparity range. Note that the HCI-Old dataset was excluded from our experimental protocol due to structural inconsistencies and incompatibility with our unified data pipeline. 

To mitigate the domain imbalance caused by the varying dataset sizes (e.g., 170 scenes in UrbanLF-Syn vs. 4 scenes in the non-Lambertian dataset), we implement a domain-balanced sampling strategy during training. To this end, we implement a domain-balanced sampling strategy to dynamically adjust the sampling probabilities, ensuring that the model is exposed to a balanced distribution of Lambertian, Mixed, and non-Lambertian domains in each mini-batch. This prevents the network from being dominated by the majority domain and facilitates the stable optimization of the domain-specific physical masks.

\textbf{Summary of Datasets.} 
Table I summarizes the key characteristics of the datasets utilized in our experiments.

\textbf{TABLE I}
\textbf{SUMMARY OF THE DATASETS USED FOR TRAINING AND EVALUATION}

\begin{table*}[!t]
\small
\caption{Dataset Domain / Scene Type Spatial Resolution.}
\label{tab:table1}
\centering
\resizebox{\textwidth}{!}{
\begin{tabular*}{\textwidth}{@{\extracolsep{\fill}}lllllll}
\toprule
Dataset \& Domain / Scene Type \& Spatial Resolution \& Angular Resolution \& Total Scenes \& Train / Val / Test Split \& Depth Source \\
\midrule
\textbf{HCI New} \& Lambertian (Complex Textures) \& $512 \times 512$ \& $9 \times 9$ \& 8 \& Scene-level partition \& Synthetic (Blender) \\
\textbf{UrbanLF-Syn} \& Mixed / Urban \& $512 \times 512$ \& $9 \times 9$ \& 170 \& Scene-level partition \& Synthetic \\
\textbf{non-Lambertian} \& non-Lambertian (Specular/Glossy) \& $512 \times 512$ \& $9 \times 9$ \& Limited (4 Train) \& Scene-level partition \& Synthetic \\
\bottomrule
\end{tabular*}
}
\end{table*}
\subsubsection{Implementation Details}

\textbf{Hardware and Software Environment.} 
All experiments were implemented using the PyTorch framework [5] and conducted on a high-performance workstation equipped with four NVIDIA RTX 3090 GPUs (24GB VRAM each). The software environment was built on CUDA 11.8 and cuDNN 8.6 to accelerate tensor operations and convolutional layers. 

\textbf{Datasets and Data Augmentation.} 
The proposed Unified Dual-Mask Physical Model was trained and evaluated on a composite dataset comprising HCInew, UrbanLF-Syn, and the non-Lambertian Dataset. The HCI-Old dataset was excluded from the experiments due to structural inconsistencies and incompatibility with our unified data pipeline. The native spatial resolution of the light field images was maintained at $512 \times 512$, with an angular resolution of $9 \times 9$. To construct the 4-directional Epipolar Plane Images (EPIs) required by the EPINet4Dir V3 [6] architecture and our dual-branch network, we extracted horizontal, vertical, and two diagonal angular slices from the sub-aperture images. 

During the training phase, a comprehensive data augmentation strategy was employed to enhance generalization and mitigate overfitting. This included random spatial cropping (extracting $256 \times 256$ patches), random horizontal and vertical flips, and photometric augmentations such as color jittering and random Gaussian noise injection. To address the severe data scarcity in the non-Lambertian domain (which contains only 4 training scenes) and to maintain exposure balance across different physical domains, a domain-balanced sampling strategy was utilized to dynamically adjust the sampling probabilities.

\textbf{Training Hyperparameters.} 
The network parameters were optimized using the AdamW optimizer [7] with a weight decay of $1 \times 10^{-4}$ to prevent overfitting. The initial learning rate was set to $1 \times 10^{-4}$ (specifically initialized at $9.99 \times 10^{-5}$) and was dynamically adjusted using a Cosine Annealing learning rate scheduler with a 5-epoch linear warmup. Extensive batch size sweeps were conducted to identify the optimal memory-performance trade-off, resulting in a batch size of 16 per GPU. With 4 GPUs, the effective batch size was 64, and gradient accumulation was not required. The model was trained for 100 epochs. The key training hyperparameters are summarized in Table II.

\textbf{TABLE II}
\textbf{SUMMARY OF KEY TRAINING HYPERPARAMETERS}

\begin{table}[!t]
\caption{Hyperparameter Value / Configuration.}
\label{tab:table2}
\centering
\begin{tabular}{ll}
\toprule
Hyperparameter \& Value / Configuration \\
\midrule
\textbf{Optimizer} \& AdamW \\
\textbf{Weight Decay} \& $1 \times 10^{-4}$ \\
\textbf{Initial Learning Rate} \& $1 \times 10^{-4}$ ($9.99 \times 10^{-5}$) \\
\textbf{LR Scheduler} \& Cosine Annealing with 5-epoch Warmup \\
\textbf{Batch Size (per GPU)} \& 16 \\
\textbf{Effective Batch Size} \& 64 \\
\textbf{Gradient Accumulation Steps} \& 1 \\
\textbf{Total Epochs} \& 100 \\
\textbf{Domain Sampling Strategy} \& Domain-balanced sampling strategy \\
\textbf{Input Spatial Resolution} \& $512 \times 512$ (cropped to $256 \times 256$) \\
\textbf{Angular Resolution} \& $9 \times 9$ (4-directional EPIs) \\
\bottomrule
\end{tabular}
\end{table}
\subsubsection{Quantitative and Qualitative Comparisons}

\textbf{Quantitative Evaluation.} 
Table III presents the quantitative depth estimation results across different domains. Compared with existing state-of-the-art methods, our proposed method achieves notable superiority in the \textbf{Overall} performance metrics. Specifically, on the HCI New dataset, it yields an overall MAE of \textbf{0.081}, which is significantly lower than the Dual-layer Net [8]. 

When evaluating the non-Lambertian branch, conventional EPI-based methods struggle with view-dependent reflections. For the baseline model, the non-Lambertian MAE significantly deteriorates to 0.585, and the Bad Pixel ratio increases substantially. In contrast, our Unified Dual-Mask Physical Model effectively mitigates these physical failures, maintaining a considerably lower error rate and demonstrating strong generalization to complex reflectance properties.

\textbf{TABLE III}
\textbf{QUANTITATIVE DEPTH ESTIMATION RESULTS ON VARIOUS DATASETS}

\begin{table}[!t]
\caption{Method HCI New (MAE $\downarrow$) UrbanLF-Syn (MAE __M....}
\label{tab:table3}
\centering
\begin{tabular*}{\linewidth}{@{\extracolsep{\fill}}lllll}
\toprule
Method \& HCI New (MAE $\downarrow$) \& UrbanLF-Syn (MAE $\downarrow$) \& non-Lambertian (MAE $\downarrow$) \& Overall (MAE $\downarrow$) \\
\midrule
\textbf{EPINet4Dir V3 [6]} \& 0.112 \& 0.145 \& 0.612 \& 0.156 \\
\textbf{Dual-layer Net [8]} \& 0.095 \& 0.121 \& 0.585 \& 0.134 \\
\textbf{Ours} \& \textbf{0.081} \& \textbf{0.098} \& \textbf{0.215} \& \textbf{0.105} \\
\bottomrule
\end{tabular*}
\end{table}
\textbf{Ablation Study.}
\paragraph{(5) Essential Role of Domain-Balanced Sampling.}
To further validate the design choices, we conduct an ablation study on the domain-balanced sampling strategy. Without this strategy, the model tends to overfit the majority domain (UrbanLF-Syn) and performs poorly on the scarce non-Lambertian scenes. By dynamically adjusting the sampling probabilities, the network successfully learns robust physical priors across all domains, confirming the necessity of our sampling mechanism for stable multi-domain optimization.

\textbf{Qualitative Analysis.}
Visual comparisons further corroborate the quantitative findings. Our model produces sharper depth boundaries and more consistent depth maps in regions with specular highlights and complex textures. In contrast, baseline methods exhibit severe blurring and structural artifacts in the non-Lambertian branch, highlighting the practical advantages of our physics-aware dual-mask design. These comprehensive empirical validations firmly establish the effectiveness of our framework and provide a natural foundation for summarizing its core contributions and overall impact.